\documentclass[10pt]{article}
\usepackage[utf8]{inputenc}
\usepackage[vietnamese]{babel}
\usepackage[paperwidth=8.5in, paperheight=11in, top=0.6in, bottom=0.6in, left=0.75in, right=0.75in]{geometry}
\usepackage{amsmath, amssymb, amsfonts}
\usepackage{graphicx}
\usepackage{xcolor}
\usepackage{tikz}
\usepackage{fancyhdr}
\usepackage{enumitem}

% --- BẢNG MÀU CHUẨN CỦA SÁCH STEWART CALCULUS ---
\definecolor{stewartcyan}{RGB}{0, 158, 224}
\definecolor{stewartred}{RGB}{195, 30, 45}
\definecolor{stewartgray}{RGB}{110, 110, 110}
\definecolor{stewartlightgray}{RGB}{240, 240, 240}

% --- THIẾT LẬP HEADER / FOOTER ---
\pagestyle{fancy}
\fancyhf{}
\fancyhead[R]{\small\textsf{\textbf{MỤC 3.7}\quad Tốc độ thay đổi trong khoa học tự nhiên và xã hội}\qquad\textbf{\textsf{225}}}
\renewcommand{\headrulewidth}{0pt}
\fancyfoot[C]{\fontsize{7}{8}\selectfont\color{stewartgray} Copyright 2021 Cengage Learning. All Rights Reserved. May not be copied, scanned, or duplicated, in whole or in part.}

\setlength{\parindent}{0pt}
\setlength{\parskip}{4pt}

\begin{document}

% =========================================================================
% PHẦN TRÊN: TIẾP TỤC CÁC BÀI TẬP CỦA MỤC TRƯỚC (BỐ CỤC 2 CỘT CÂN BẰNG)
% =========================================================================
\noindent
\begin{minipage}[t]{0.475\textwidth}
    \begin{enumerate}[label=\textbf{\arabic*.}, leftmargin=*, itemsep=2.5pt, topsep=0pt]
        \setcounter{enumi}{72}
        \item $h(t) = \cot^{-1}(t) + \cot^{-1}(1/t)$
        \item $R(t) = \arcsin(1/t)$
        \item $y = x\sin^{-1}x + \sqrt{1 - x^2}$
        \item $y = \cos^{-1}(\sin^{-1}t)$
        \item $y = \tan^{-1}\!\left(\dfrac{x}{a}\right) + \ln\sqrt{\dfrac{x - a}{x + a}}$
        \item $y = \arctan\sqrt{\dfrac{1 - x}{1 + x}}$
    \end{enumerate}

    \vspace{2pt}
    {\color{stewartcyan}\rule{\linewidth}{0.6pt}}
    \vspace{1pt}

    \noindent
    \tikz[baseline=-0.6ex]{
        \draw[stewartcyan, fill=white, line width=0.8pt] (0,0) rectangle (0.32, 0.30);
        \draw[stewartcyan, line width=0.8pt] (0.04, 0.08) .. controls (0.12, 0.26) and (0.20, 0.05) .. (0.28, 0.22);
    }\hspace{4pt}%
    \textbf{\color{stewartcyan}79--80}\; \small Tìm $f'(x)$. Kiểm tra xem câu trả lời của bạn có hợp lý hay không bằng cách so sánh đồ thị của $f$ và $f'$.
    
    \vspace{3pt}
    \noindent
    \textbf{79.}\; $f(x) = \sqrt{1 - x^2}\;\arcsin x$\hfill
    \textbf{80.}\; $f(x) = \arctan(x^2 - x)$

    \vspace{1pt}
    {\color{stewartcyan}\rule{\linewidth}{0.6pt}}
    \vspace{3pt}

    \noindent
    \textbf{81.}\; Chứng minh công thức tính $(d/dx)(\cos^{-1}x)$ bằng phương pháp tương tự như đối với $(d/dx)(\sin^{-1}x)$.

    \vspace{4pt}
    \noindent
    \textbf{82.}\; (a) Một cách định nghĩa $\sec^{-1}x$ là coi $y = \sec^{-1}x \iff \sec y = x$ và $0 \le y < \pi/2$ hoặc $\pi \le y < 3\pi/2$. Hãy chứng minh rằng, với định nghĩa này,
    \[
    \frac{d}{dx}(\sec^{-1}x) = \frac{1}{x\sqrt{x^2 - 1}}
    \]
\end{minipage}%
\hfill
\begin{minipage}[t]{0.485\textwidth}
    \noindent
    (b) Một cách định nghĩa $\sec^{-1}x$ khác đôi khi cũng được sử dụng là coi $y = \sec^{-1}x \iff \sec y = x$ và $0 \le y \le \pi,\; y \ne \pi/2$. Hãy chứng minh rằng, với định nghĩa này,
    \[
    \frac{d}{dx}(\sec^{-1}x) = \frac{1}{|x|\sqrt{x^2 - 1}}
    \]

    \noindent
    \textbf{83. Đạo hàm của hàm ngược}\quad Giả sử $f$ là hàm đơn ánh khả vi và hàm ngược $f^{-1}$ của nó cũng khả vi. Sử dụng phương pháp đạo hàm ẩn để chứng minh rằng
    \[
    (f^{-1})'(x) = \frac{1}{f'(f^{-1}(x))}
    \]
    với điều kiện mẫu số khác 0.

    \vspace{2pt}
    \noindent
    \textbf{\color{stewartcyan}84--86}\; Sử dụng công thức trong Bài tập 83.
    \begin{enumerate}[label=\textbf{\arabic*.}, leftmargin=*, itemsep=1.5pt, topsep=2pt]
        \setcounter{enumi}{83}
        \item Nếu $f(4) = 5$ và $f'(4) = \frac{2}{3}$, hãy tìm $(f^{-1})'(5)$.
        \item Nếu $f(x) = x + e^x$, hãy tìm $(f^{-1})'(1)$.
        \item Nếu $f(x) = x^3 + 3\sin x + 2\cos x$, hãy tìm $(f^{-1})'(2)$.
    \end{enumerate}

    \vspace{2pt}
    {\color{stewartcyan}\rule{\linewidth}{0.6pt}}
    \vspace{3pt}

    \noindent
    \textbf{87.}\; Giả sử rằng $f$ và $g$ là các hàm khả vi và đặt $h(x) = f(x)^{g(x)}$. Sử dụng đạo hàm logarit để thiết lập công thức
    \[
    h' = g \cdot f^{g-1} \cdot f' + (\ln f) \cdot f^g \cdot g'
    \]

    \noindent
    \textbf{88.}\; Sử dụng công thức trong Bài tập 87 để tìm đạo hàm.
    
    \vspace{2pt}
    \noindent
    (a) $h(x) = x^3$\hfill (b) $h(x) = 3^x$\hfill (c) $h(x) = (\sin x)^x$
\end{minipage}

\vspace{18pt}

% =========================================================================
% DẢI TIÊU ĐỀ MỤC 3.7
% =========================================================================
\noindent
\begin{minipage}[c]{0.26\textwidth}
    \color{stewartcyan}
    \rule{\linewidth}{1.2pt}\\[3.5pt]
    \rule{\linewidth}{1.2pt}\\[3.5pt]
    \rule{\linewidth}{1.2pt}\\[3.5pt]
    \rule{\linewidth}{1.2pt}
\end{minipage}%
\hspace{8pt}%
\begin{minipage}[c]{0.70\textwidth}
    \noindent\textsf{\textbf{\fontsize{24}{26}\selectfont\color{stewartcyan}3.7}}%
    \hspace{6pt}%
    {\color{stewartcyan}\rule[-4pt]{1.5pt}{24pt}}%
    \hspace{8pt}%
    \textsf{\textbf{\fontsize{13.5}{16}\selectfont Tốc độ thay đổi trong Khoa học Tự nhiên và Xã hội}}
\end{minipage}

\vspace{10pt}

% =========================================================================
% PHẦN DƯỚI: NỘI DUNG LÝ THUYẾT VÀ HÌNH MINH HỌA BÊN LỀ
% =========================================================================
\noindent
\begin{minipage}[t]{0.26\textwidth}
    % Khoảng cách dóng thẳng hàng hình minh họa với đoạn nhắc đến "Hình 1"
    \vspace{70pt}
    
    \centering
    \includegraphics[width=\linewidth]{images/p0260_fig1.png}
    
    \vspace{6pt}
    {\small\raggedright
    $m_{PQ} = \text{tốc độ thay đổi trung bình}$\\
    $m = f'(x_1) = \text{tốc độ thay đổi tức thời}$\par
    }
    
    \vspace{6pt}
    \noindent{\bfseries\textsf{HÌNH 1}}
\end{minipage}%
\hfill
\begin{minipage}[t]{0.71\textwidth}
    Chúng ta biết rằng nếu $y = f(x)$, thì đạo hàm $dy/dx$ có thể được hiểu là tốc độ thay đổi của $y$ theo $x$. Trong mục này, chúng ta sẽ xem xét một số ứng dụng của ý tưởng này vào vật lý, hóa học, sinh học, kinh tế học và các ngành khoa học khác.

    Hãy nhớ lại từ Mục 2.7 ý tưởng cơ bản đằng sau tốc độ thay đổi. Nếu $x$ thay đổi từ $x_1$ đến $x_2$, thì độ biến thiên của $x$ là
    \[
    \Delta x = x_2 - x_1
    \]
    và độ biến thiên tương ứng của $y$ là
    \[
    \Delta y = f(x_2) - f(x_1)
    \]
    Tỉ số biến thiên
    \[
    \frac{\Delta y}{\Delta x} = \frac{f(x_2) - f(x_1)}{x_2 - x_1}
    \]
    là \textbf{tốc độ thay đổi trung bình của $y$ theo $x$} trên đoạn $[x_1, x_2]$ và có thể được hiểu là hệ số góc của cát tuyến $PQ$ trong Hình 1. Giới hạn của tỉ số này khi $\Delta x \to 0$ là đạo hàm $f'(x_1)$, do đó có thể được hiểu là \textbf{tốc độ thay đổi tức thời của $y$ theo $x$} hoặc hệ số góc của tiếp tuyến tại điểm $P(x_1, f(x_1))$. Sử dụng ký hiệu Leibniz, chúng ta viết quá trình này dưới dạng
    \[
    \frac{dy}{dx} = \lim_{\Delta x \to 0} \frac{\Delta y}{\Delta x}
    \]
\end{minipage}

\end{document}